% !TEX root = ../../main.tex
\section{الاستنتاجات} \label{part1:sec5}
% \thispagestyle{empty}

من خلال هذه الدراسة والمحاكاة العملية باستخدام برنامج \el{\textbf{MATLAB R2023b}}، يمكن تلخيص الاستنتاجات التالية:

\begin{enumerate}
    \item \textbf{فعالية التصميم:} أثبتت مرشحات التردد الرنيني \el{(\textbf{Tuned Filters})} كفاءة عالية في امتصاص تيارات التوافقيات عند ترددات الرنين المحددة، مما أدى لخفض نسبة \el{$THD$} بشكل تدريجي مع كل إضافة.
    \item \textbf{تحسين جودة القدرة:} ساهمت المرشحات في تحسين معامل الاستطاعة وتقريب شكل موجة التيار والجهد من الشكل الجيبي النقي، مما يقلل من الفواقد والحرارة في عناصر الشبكة.
    \item \textbf{مطابقة المعايير:} بفضل إضافة مرشحات الرتب 7 و 11، تمكنا من الحفاظ على جودة الجهد ضمن الحدود العالمية المسموح بها (غالباً أقل من 3\% أو 5\% حسب المواصفة).
    \item \textbf{أهمية النمذجة:} أكدت الدراسة أن ضبط إعدادات المحاكاة (مثل وضع \el{\textbf{Discrete}} وزمن العينة الدقيق) واستخدام الإزاحة الطورية الصحيحة ($120^\circ$) هي عوامل حاسمة للحصول على نتائج مطابقة للواقع الفيزيائي.
\end{enumerate}
